\todo{Lecture 22}
\chapter{Magnétostatique}
Situations avec courants constants dans le temps ("écoulement stationnaire des charges").
\paragraph{Observations expérimentales}
\begin{itemize}
\item Boussole (aiguille aimantée) est déviée par un aimant.
\item Boussole est déviée par un fil portant un courant.
\item Aimant exerce force sur fil portant un courant.
\item Fil portant un courant exerce force sur un autre fil portant un courant. N'est pas une interaction électrostatique !
\item Aimant pas charge : pas d'effet sur un diélectrique (p.ex. bouts de papier)
\item Aimant influence boussole a travers une cage de Faraday.
\item Les fils portant un courant ne sont pas charges.
\end{itemize}

\begin{definition}
Le champ magnétique $\boldsymbol{B}$ est génère par un aimant ou courant, et exerce une force sur un aimant ou courant. La direction de $\boldsymbol{B}$ est défini par l'orientation d'une boussole.
\end{definition}
\section{Definition du champ magnétique et force de Lorentz}

Le champ magnétique $\boldsymbol{B}$ exerce une force sur un fil portant un courant. La force est exercée sur les porteurs de charge en mouvement.

Des experiences montrent que la force exercee sur une charge $q$ de vitesse $\boldsymbol{v}$ est
\begin{itemize}
\item perpendiculaire a $\boldsymbol{v}$ et $\boldsymbol{B}$,
\item proportionel a $\lVert \boldsymbol{v} \rVert$ et $q$.
\end{itemize}
donc 
$$
\boldsymbol{F}=q(\boldsymbol{x}\land \boldsymbol{B}).
$$
Le champ magnétique $\boldsymbol{B}$ se mesure en tesla $\mathrm{T}=\mathrm{NC}^{-1}=\mathrm{Sm}^{-1}=\mathrm{kgA}^{-1}\mathrm{s}^{-2}$.
\begin{definition}[Force de Lorentz]
En presence de $\boldsymbol{E}$ la force est
$$
\boldsymbol{F}=q(\boldsymbol{E}+\boldsymbol{v}\land \boldsymbol{B}).
$$
\end{definition}
\paragraph{Quelques valeurs typiques de $\boldsymbol{B}$}
\begin{itemize}
\item Le champ terrestre $\approx 0.5\cdot10^{-4}\,\mathrm{T}$
\item Le champ maximale dans un supraconducteur $\le 30\,\mathrm{T}$.
\end{itemize}

La force d'un champ $\boldsymbol{B}$ externe sur un element de fil $d\boldsymbol{\ell }$ portant un courant $I$ :
$$
d\boldsymbol{F}=nS\,d \ell\, q(\boldsymbol{v}\land \boldsymbol{B}).
$$
Comme $\boldsymbol{v}\parallel d\boldsymbol{\ell}$, alors $d \ell\,\boldsymbol{v}=\boldsymbol{v}\,d\boldsymbol{\ell}$ et
$$
d\boldsymbol{F}=nSvq(d\boldsymbol{\ell}\land \boldsymbol{B})=I\,d\boldsymbol{\ell}\land \boldsymbol{B}.
$$
